% Generated from the original tree by tmp/gen_mindmaps.py; edit the text here directly.
\mapfigure{%
\node[mmroot] (r) {What stacking costs, and AMD's answers};
\node[mmbox,mmlevel1,fill=mmF4AAA2] (b0) at (0.00,3.60) {Heat path};
\draw[mmbig] (r) -- (b0);
\node[mmlevel2,mmbox,fill=mmCFE7E1] (b0c0) at (-3.05,5.35) {Extra layers raise $R$};
\draw[mmsmall] (b0.north) to[out=90,in=-90] (b0c0.south);
\node[mmlevel2,mmbox,fill=mmCFE7E1] (b0c1) at (0.00,5.35) {Leakage feedback};
\draw[mmsmall] (b0.north) to[out=90,in=-90] (b0c1.south);
\node[mmlevel2,mmbox,fill=mmCFE7E1] (b0c2) at (3.05,5.35) {Hot spot: $\Delta T = RP$};
\draw[mmsmall] (b0.north) to[out=90,in=-90] (b0c2.south);
\node[mmbox,mmlevel1,fill=mmF9D7AF] (b1) at (4.47,1.36) {Voltage and clock};
\draw[mmbig] (r) -- (b1);
\node[mmlevel2,mmbox,fill=mmDCCFF6] (b1c0) at (8.07,2.81) {1.2\,V limit (7800X3D)};
\draw[mmsmall] (b1.east) to[out=0,in=180] (b1c0.west);
\node[mmlevel2,mmbox,fill=mmDCCFF6] (b1c1) at (8.07,1.36) {4.8 against 5.15\,GHz (7950X3D)};
\draw[mmsmall] (b1.east) to[out=0,in=180] (b1c1.west);
\node[mmlevel2,mmbox,fill=mmDCCFF6] (b1c2) at (8.07,-0.09) {Lower all-core boost (9684X)};
\draw[mmsmall] (b1.east) to[out=0,in=180] (b1c2.west);
\node[mmbox,mmlevel1,fill=mmCFE0F4] (b2) at (2.76,-3.56) {Power delivery};
\draw[mmbig] (r) -- (b2);
\node[mmlevel2,mmbox,fill=mmF9D7AF] (b2c0) at (6.36,-2.83) {RVDD and VDDM};
\draw[mmsmall] (b2.east) to[out=0,in=180] (b2c0.west);
\node[mmlevel2,mmbox,fill=mmF9D7AF] (b2c1) at (6.36,-4.28) {Power TSVs take area};
\draw[mmsmall] (b2.east) to[out=0,in=180] (b2c1.west);
\node[mmoval,mmlevel1,fill=mmDCCFF6] (b3) at (-2.76,-3.56) {Manufacturing};
\draw[mmbig] (r) -- (b3);
\node[mmlevel2,mmbox,fill=mmF4C2BB] (b3c0) at (-6.36,-2.11) {Thinning to 20\,\textmu m};
\draw[mmsmall] (b3.west) to[out=180,in=0] (b3c0.east);
\node[mmlevel2,mmbox,fill=mmF4C2BB] (b3c1) at (-6.36,-3.56) {Known-good dies};
\draw[mmsmall] (b3.west) to[out=180,in=0] (b3c1.east);
\node[mmlevel2,mmbox,fill=mmF4C2BB] (b3c2) at (-6.36,-5.01) {Bonding yield};
\draw[mmsmall] (b3.west) to[out=180,in=0] (b3c2.east);
\node[mmbox,mmlevel1,fill=mmD6EBC6] (b4) at (-4.47,1.36) {AMD's answers};
\draw[mmbig] (r) -- (b4);
\node[mmlevel2,mmbox,fill=mmDCCFF6] (b4c0) at (-8.07,2.81) {Cache over the L3 only};
\draw[mmsmall] (b4.west) to[out=180,in=0] (b4c0.east);
\node[mmlevel2,mmbox,fill=mmDCCFF6] (b4c1) at (-8.07,1.36) {Structural silicon over cores};
\draw[mmsmall] (b4.west) to[out=180,in=0] (b4c1.east);
\node[mmlevel2,mmbox,fill=mmDCCFF6] (b4c2) at (-8.07,-0.09) {Zen~5: cache below cores};
\draw[mmsmall] (b4.west) to[out=180,in=0] (b4c2.east);
}{[Overview of the costs of stacking and AMD's responses]{Overview of the costs of stacking discussed in \cref{sec:thermal}, with the measurements cited there, and of AMD's design responses. The sources attribute the voltage and clock limits at least partly to the stacked cache's voltage tolerance as well as to heat (\cref{sec:thermal-evidence}).}}{fig:mindmap-stack-costs}
